\section*{الفصل \ref{ch:g6:propdata} --- التناسب والمعطيات}

\begin{solution}{exo:g6:propdata:1}
التفاح: \omterm{def:g3:division:remainder}{خوارج القسمة} $\frac52 = 2.5$ و $\frac{7.5}{3} = 2.5$ و
$\frac{12.5}{5} = 2.5$ --- وكلها متساوية: إذن متناسب (والثمن
$2.50$ للكيلوغرام).

السنّ والطول: $\frac{85}{2} = 42.5$ لكن $\frac{103}{4} = 25.75$ --- غير
متساويين: إذن غير \omterm{def:g6:propdata:def}{متناسبين}.
\end{solution}

\begin{solution}{exo:g6:propdata:2}
الكرواسان الواحد: $4.80 \div 4 = 1.20$ يورو. وسبعة:
$7 \times 1.20 = 8.40$ يورو.
\end{solution}

\begin{solution}{exo:g6:propdata:3}
المعامل: $18 \div 10 = 1.8$ يورو للتر الواحد. ثم
$25$ لترًا ثمنها $25 \times 1.8 = 45$ يورو؛ و $72$ يورو تشتري
$72 \div 1.8 = 40$ لترًا؛ و $60$ لترًا ثمنها $60 \times 1.8 = 108$ يورو.
\end{solution}

\begin{solution}{exo:g6:propdata:4}
من أجل $250$ كم: أي $2.5$ أمثال وقود $100$ كم، إذن
$2.5 \times 6 = 15$ لترًا. ومن أجل $350$ كم: $3.5 \times 6 = 21$ لترًا. وبمقدار
$27$ لترًا: $27 \div 6 = 4.5$ مئات من الكيلومترات، أي\ $450$ كم.
\end{solution}

\begin{solution}{exo:g6:propdata:5}
$50\,\%$ من $84$: النصف، أي $42$. \quad
$25\,\%$ من $200$: الرُّبع، أي $50$. \quad
$10\,\%$ من $63$: العُشر، أي $6.3$. \quad
$30\,\%$ من $70$: $3 \times 7 = 21$.
\end{solution}

\begin{solution}{exo:g6:propdata:6}
$60\,\%$ من $450$: $\frac{60}{100} \times 450 = 270$ تلميذًا يأكلون في
المطعم؛ و $450 - 270 = 180$ لا يأكلون فيه.
\end{solution}

\begin{solution}{exo:g6:propdata:7}
أقلّ من $10$ كتب: الثلاثاء ($8$) والخميس ($6$). والأربعاء
ناقص الخميس: $17 - 6 = 11$ كتابًا أكثر.
\end{solution}

\begin{solution}{exo:g6:propdata:8}
مخطط أعمدة بخمسة أعمدة اِرتفاعاتها $14$ و $16$ و $13$ و $17$ و $20$
(والتدريج كل $2$ أو $5$ درجات يعمل جيدًا). وأدفأ يوم: الجمعة
($20$ درجة).
\end{solution}

\begin{solution}{exo:g6:propdata:9}
من أجل $10$ أشخاص، اِضرب في $\frac{10}{6} = \frac53$: فالدقيق
$450 \times \frac{10}{6} = 750$ غ. والبيض:
$3 \times \frac{10}{6} = 5$ --- ولحسن الحظ عدد صحيح. (ولو لم يكن
كذلك، لقُرِّب \emph{بالزيادة} حتى يكفي.)
\end{solution}

\begin{solution}{exo:g6:propdata:10}
\emph{1.} في $2$ ساعة: $48$ كم. وفي نصف ساعة: $12$ كم. وفي
$1$ ساعة و $30$: $24 + 12 = 36$ كم.

\emph{2.} $60 \div 24 = 2.5$ ساعة، أي\ $2$ ساعة و $30$ دقيقة.
\end{solution}

\begin{solution}{exo:g6:propdata:11}
الكرة: التخفيض $20\,\%$ من $15 = 3$ يورو، والثمن الجديد $12$ يورو.
والمضرب: التخفيض $8$ يورو، والثمن الجديد $32$ يورو. والثمن الجديد هو
$80\,\%$ من القديم في كل حالة: إذن متناسب، ومعامله
$0.8$.
\end{solution}

\begin{solution}{exo:g6:propdata:12}
بعد $t$ ساعة، تكون الشمعة الغليظة قد احترقت $\frac t6$ من اِرتفاعها:
فتقف عند $1 - \frac t6$ من الارتفاع الابتدائي؛ والرفيعة عند
$1 - \frac t4$. ونريد
\[
1 - \frac t4 = \frac12 \left(1 - \frac t6\right)
\quad\text{أي}\quad
1 - \frac t4 = \frac12 - \frac t{12}.
\]
ومنه $\frac12 = \frac t4 - \frac t{12} = \frac{3t - t}{12} =
\frac{t}{6}$، إذن $t = 3$ ساعات. والتحقّق: بعد $3$ ساعات تقف الشمعة الغليظة
عند $1 - \frac36 = \frac12$ والرفيعة عند
$1 - \frac34 = \frac14$ --- وهو فعلًا نصفه.
\end{solution}

\begin{solution}{pb:g6:propdata:1}
\textbf{1.} $100\,000$ سم $= 1\,000$ م $= 1$ كم: فالسنتيمتر الواحد
على الخريطة كيلومتر واحد على الأرض.

\textbf{2.} $7.5$ سم $\rightarrow 7.5$ كم بين القريتين؛
و $12$ سم $\rightarrow 12$ كم من شاطئ البحيرة.

\textbf{3.} $20$ كم $\rightarrow 20$ سم على الخريطة.

\textbf{4.} $5$ كم $= 500\,000$ سم، إذن السلّم
$1:500\,000$. وخريطة المشي ($1:100\,000$) هي الأكثر
تفصيلًا: فهي تستعمل $5$ سم من الورق حيث لا تنفق الأخرى إلا
$1$ سم.

\textbf{5.}
\begin{center}
\renewcommand{\arraystretch}{1.2}
\begin{tabular}{l|ccc}
الخريطة (سم) & $1$ & $2$ & $7.5$ \\
\hline
الحقيقة (كم) & $1$ & $2$ & $7.5$ \\
\end{tabular}
\end{center}
المسافة الحقيقية $=$ المسافة على الخريطة $\times 1$ (بهذه الوحدات):
أي المضروب فيه \emph{نفسه} دائمًا، وهو بالضبط
تعريف \omterm{def:g6:propdata:def}{التناسب}
(\cref{def:g6:propdata:def}). والمعامل هنا $1$
كيلومتر لكل سنتيمتر.

\textbf{6.} $4$ م $= 400$ سم و $400 \div 50 = 8$؛
و $3$ م $= 300$ سم و $300 \div 50 = 6$: فالتصميم يبيّن مستطيلًا
$8$ سم $\times$ $6$ سم.

\textbf{7.} السرير: $200 \div 50 = 4$ سم في $90 \div 50 = 1.8$ سم.
والباب: $80 \div 50 = 1.6$ سم.

\textbf{8.} \omterm{def:g6:measure:area}{المساحة} الحقيقية: $4 \times 3 = 12$ م$^2$. \omterm{def:g6:measure:area}{ومساحة} التصميم:
$8 \times 6 = 48$ سم$^2$. وبالتحويل:
$12$ م$^2 = 12 \times 10\,000 = 120\,000$ سم$^2$، ثم
\[
120\,000 \div 48 = 2\,500 .
\]
\omterm{def:g3:division:remainder}{وخارج القسمة} ليس $50$ بل $2\,500 = 50 \times 50$.

\textbf{9.} كل سم$^2$ من الورق يمثّل مربعًا حقيقيًّا
$50$ سم في $50$ سم، وهو يحوي
$50 \times 50 = 2\,500$ سم$^2$. إذن حين تُقسم الأطوال على
$50$، تُقسم المساحات على $50 \times 50 = 2\,500$: فالمساحات تتبع
\emph{مربّع} السلّم.

\textbf{10.} $6 \times 2\,500 = 15\,000$ سم$^2$، ثم
$15\,000 \div 10\,000 = 1.5$ م$^2$.

\textbf{11.} توحي الصورة بأن الأحد باع \emph{ضعف} ما
باعه السبت (فعموده ضعف الارتفاع). والحقيقة: أن الزيادة $5$ يورو من
أصل $105$، و $5 \div 105 \approx 0.048$: أي نحو $5\,\%$ أكثر،
لا $100\,\%$ أكثر. وقد أهمل المدير التحذير بالتحقّق
من أن المحور يبدأ من $0$
(\cref{met:g6:propdata:read}).

\textbf{12.} بمحور يبدأ من $0$، يرتفع العمودان إلى $105$ و
$110$ وحدة صغيرة: أي عمودان بالارتفاع نفسه تقريبًا، وهو
الصورة الأمينة \omterm{ex:g1:subtraction:difference}{لفرق} $5\,\%$.

\textbf{13.} صعودًا: الزيادة $20$ طابعًا من بداية
$40$: أي $\frac{20}{40} = 50\,\%$. ونزولًا: النقصان $20$
طابعًا من بداية $60$: أي $\frac{20}{60} = \frac13 \approx
33\,\%$. الطوابع $20$ نفسها، وقيمتان اِبتدائيتان مختلفتان --- فالنسبة
المئوية دائمًا \omterm{def:g6:fractions:def}{كسر} \emph{من شيء ما}، وقد تغيّر
«الشيء».

\textbf{14.} بعد تخفيض $20\,\%$: $100 - 20 = 80$ يورو.
وتنطبق $10\,\%$ الإضافية على $80$: فالتخفيض $8$ يورو، والثمن النهائي
$72$ يورو. \omterm{def:g1:addition:def}{ومجموع} التخفيض: $28$ يورو من أصل $100$، أي
$28\,\%$ --- لا $30\,\%$. فقد عمل التخفيض الثاني على
الثمن المخفَّض أصلًا، فيوروهاته أصغر: فالنسب المئوية من
مقادير مختلفة لا تُجمع.

\textbf{15.} على تلك الخريطة، يمثّل $1$ سم مسافة $5$ كم، إذن $1$
سم$^2$ يمثّل $5 \times 5 = 25$ كم$^2$ (قاعدة السؤال 9).
والغابة: $3 \times 25 = 75$ كم$^2$. فالأطوال تتبع
السلّم؛ والمساحات تتبع \emph{مربّعه}.
\end{solution}
